% ══════════════════════════════════════════════════════════════════
% RÉSUMÉ (français) + ABSTRACT (anglais)
% ══════════════════════════════════════════════════════════════════

\chapter*{Résumé}
\addcontentsline{toc}{chapter}{Résumé}
\markboth{Résumé}{Résumé}
\thispagestyle{plain}

L'agriculture tchadienne, pilier de l'économie nationale, souffre
d'un déficit d'outils d'aide à la décision adaptés aux réalités
locales. Ce mémoire présente la conception et l'implémentation
d'un cadre décisionnel basé sur l'Intelligence Artificielle pour
l'optimisation du suivi parcellaire agricole au Tchad, dénommé
\textbf{Agro-IA Tchad}.

\medskip

Le système intègre quatre modules complémentaires. Le
\textbf{module Irrigation} repose sur l'algorithme Random Forest,
entraîné sur 16~435 observations climatiques réelles issues de la
base NASA POWER pour cinq villes tchadiennes (2015--2023) ; il
atteint une accuracy de 78,51\,\% et un AUC-ROC de 0,8782. Le
\textbf{module Fertilisation} utilise XGBoost sur 1~599
observations et recommande l'engrais optimal avec une accuracy de
95,42\,\% (F1-score macro de 95,47\,\%). Le \textbf{module
	Détection des maladies} combine un réseau de neurones convolutif
MobileNetV2 --- entraîné par transfert d'apprentissage sur un
dataset local de 287 images de cultures tchadiennes (accuracy de
52,27\,\%) --- avec l'analyse d'indices de végétation NDVI issus
du satellite Sentinel-2 ou estimés par caméra smartphone (indice
VARI). Enfin, le \textbf{module Historique} assure la traçabilité
de toutes les analyses grâce à une base de données SQLite avec
export CSV.

\medskip

L'application web, développée avec le framework Streamlit et
intégrant l'API météorologique Open-Meteo en temps réel, est
déployée sur Streamlit Cloud et accessible depuis tout navigateur.

\medskip

\noindent\textbf{Mots-clés :} Intelligence Artificielle, Random
Forest, XGBoost, CNN MobileNetV2, NDVI, Agriculture de précision,
Tchad, Streamlit, SQLite.

\newpage

% ──────────────────────────────────────────────────────────────────
\chapter*{Abstract}
\addcontentsline{toc}{chapter}{Abstract}
\markboth{Abstract}{Abstract}
\thispagestyle{plain}

Chadian agriculture, a pillar of the national economy, suffers
from a lack of decision-support tools adapted to local realities.
This thesis presents the design and implementation of an
Artificial Intelligence-based decision framework for optimizing
agricultural plot monitoring in Chad, named \textbf{Agro-IA
	Tchad}.

\medskip

The system integrates four complementary modules. The
\textbf{Irrigation module} is based on the Random Forest
algorithm, trained on 16,435 real climate observations from the
NASA POWER database for five Chadian cities (2015--2023),
achieving 78.51\,\% accuracy and an AUC-ROC of 0.8782. The
\textbf{Fertilization module} uses XGBoost on 1,599 observations
and recommends the optimal fertilizer with 95.42\,\% accuracy
(macro F1-score of 95.47\,\%). The \textbf{Disease Detection
	module} combines a MobileNetV2 convolutional neural network ---
trained by transfer learning on a local dataset of 287 images of
Chadian crops (52.27\,\% accuracy) --- with NDVI vegetation
indices from the Sentinel-2 satellite or estimated via smartphone
camera (VARI index). Finally, the \textbf{History module} ensures
traceability of all analyses through a SQLite database with CSV
export.

\medskip

The web application, developed with the Streamlit framework and
integrating the Open-Meteo real-time weather API, is deployed on
Streamlit Cloud and accessible from any browser.

\medskip

\noindent\textbf{Keywords:} Artificial Intelligence, Random
Forest, XGBoost, CNN MobileNetV2, NDVI, Precision Agriculture,
Chad, Streamlit, SQLite.